% Options for packages loaded elsewhere
\ifLuaTeX
\usepackage[bidi=basic]{babel}
\else
\usepackage[bidi=default]{babel}
\fi
% get rid of language-specific shorthands (see #6817):
\let\LanguageShortHands\languageshorthands
\def\languageshorthands#1{}


\setlength{\emergencystretch}{3em} % prevent overfull lines

\providecommand{\tightlist}{%
  \setlength{\itemsep}{0pt}\setlength{\parskip}{0pt}}



 


\usepackage{multirow}
\usepackage{longtable}
\makeatletter
\@ifpackageloaded{tcolorbox}{}{\usepackage[skins,breakable]{tcolorbox}}
\@ifpackageloaded{fontawesome5}{}{\usepackage{fontawesome5}}
\definecolor{quarto-callout-color}{HTML}{909090}
\definecolor{quarto-callout-note-color}{HTML}{0758E5}
\definecolor{quarto-callout-important-color}{HTML}{CC1914}
\definecolor{quarto-callout-warning-color}{HTML}{EB9113}
\definecolor{quarto-callout-tip-color}{HTML}{00A047}
\definecolor{quarto-callout-caution-color}{HTML}{FC5300}
\definecolor{quarto-callout-color-frame}{HTML}{acacac}
\definecolor{quarto-callout-note-color-frame}{HTML}{4582ec}
\definecolor{quarto-callout-important-color-frame}{HTML}{d9534f}
\definecolor{quarto-callout-warning-color-frame}{HTML}{f0ad4e}
\definecolor{quarto-callout-tip-color-frame}{HTML}{02b875}
\definecolor{quarto-callout-caution-color-frame}{HTML}{fd7e14}
\makeatother
\makeatletter
\@ifpackageloaded{bookmark}{}{\usepackage{bookmark}}
\makeatother
\makeatletter
\@ifpackageloaded{caption}{}{\usepackage{caption}}
\AtBeginDocument{%
\ifdefined\contentsname
  \renewcommand*\contentsname{Table des matières}
\else
  \newcommand\contentsname{Table des matières}
\fi
\ifdefined\listfigurename
  \renewcommand*\listfigurename{Liste des Figures}
\else
  \newcommand\listfigurename{Liste des Figures}
\fi
\ifdefined\listtablename
  \renewcommand*\listtablename{Liste des Tables}
\else
  \newcommand\listtablename{Liste des Tables}
\fi
\ifdefined\figurename
  \renewcommand*\figurename{Figure}
\else
  \newcommand\figurename{Figure}
\fi
\ifdefined\tablename
  \renewcommand*\tablename{Table}
\else
  \newcommand\tablename{Table}
\fi
}
\@ifpackageloaded{float}{}{\usepackage{float}}
\floatstyle{ruled}
\@ifundefined{c@chapter}{\newfloat{codelisting}{h}{lop}}{\newfloat{codelisting}{h}{lop}[chapter]}
\floatname{codelisting}{Listing}
\newcommand*\listoflistings{\listof{codelisting}{Liste des Listings}}
\makeatother
\makeatletter
\makeatother
\makeatletter
\@ifpackageloaded{caption}{}{\usepackage{caption}}
\@ifpackageloaded{subcaption}{}{\usepackage{subcaption}}
\makeatother
\usepackage{bookmark}
\IfFileExists{xurl.sty}{\usepackage{xurl}}{} % add URL line breaks if available
\urlstyle{same}
\hypersetup{
  pdftitle={Campagne 2025 des Observatoires ruraux d'Alaotra et Marovoay},
  pdfauthor={Veromanitra Ramizason, Bezaka Rivolala, Thierry Razanakoto, Holimalala Randriamanampisoa, Florent Bédécarrats},
  pdflang={fr},
  colorlinks=true,
  linkcolor={blue},
  filecolor={Maroon},
  citecolor={Blue},
  urlcolor={Blue},
  pdfcreator={LaTeX via pandoc}}


\title{Campagne 2025 des Observatoires ruraux d'Alaotra et Marovoay}
\usepackage{etoolbox}
\makeatletter
\providecommand{\subtitle}[1]{% add subtitle to \maketitle
  \apptocmd{\@title}{\par {\large #1 \par}}{}{}
}
\makeatother
\subtitle{Rapport consolidé d'enquêtes}
\author{Veromanitra Ramizason, Bezaka Rivolala, Thierry Razanakoto,
Holimalala Randriamanampisoa, Florent Bédécarrats}
\date{2026-03-15}
\begin{document}

\begin{table}

\caption{\label{tbl-menaces}Menaces perçues sur l'environnement et
responsables identifiés (\% de Oui)}

\centering{

\fontsize{12.0pt}{14.0pt}\selectfont
\begin{tabular*}{\linewidth}{@{\extracolsep{\fill}}lrr}
\toprule
 & {\bfseries {Alaotra}} & {\bfseries {Marovoay}} \\ 
\midrule\addlinespace[2.5pt]
\multicolumn{3}{l}{{\bfseries {Perception globale}}} \\[2.5pt] 
\midrule\addlinespace[2.5pt]
Considèrent que l'environnement est menacé & 82 % & 79 % \\ 
\midrule\addlinespace[2.5pt]
\multicolumn{3}{l}{{\bfseries {Principales menaces (parmi ceux qui déclarent menacé)}}} \\[2.5pt] 
\midrule\addlinespace[2.5pt]
Feux & 84 % & 95 % \\ 
Défrichage & 56 % & 52 % \\ 
Prélèvement (faune) & 40 % & 20 % \\ 
Prélèvement (flore) & 40 % & 21 % \\ 
Prospection minière & 29 % & 6 % \\ 
Nomadisme (passage non-résidents) & 16 % & 16 % \\ 
Zébu / pâturage & 20 % & 10 % \\ 
Tourisme & 2 % & 3 % \\ 
\midrule\addlinespace[2.5pt]
\multicolumn{3}{l}{{\bfseries {Principaux responsables}}} \\[2.5pt] 
\midrule\addlinespace[2.5pt]
L'ensemble des riverains & 87 % & 70 % \\ 
Les riverains migrants & 47 % & 37 % \\ 
L'État (Fanjakana) & 10 % & 8 % \\ 
Les braconniers & 30 % & 10 % \\ 
Le gestionnaire & 9 % & 3 % \\ 
Les nomades ou semi-nomades & 22 % & 17 % \\ 
Les ONG ou entreprises & 3 % & 4 % \\ 
Les chercheurs & 0 % & 2 % \\ 
Les touristes & 1 % & 2 % \\ 
\bottomrule
\end{tabular*}
\begin{minipage}{\linewidth}
\textbf{Source : Enquête auprès des OR 2025}\\
\end{minipage}

}

\end{table}%
\end{document}
